\chapter{Propensity Scores and Weighting}
\label{chap:propensity}

\section{The balancing property}
\label{sec:ps-balancing}

\begin{definition}[Propensity score]
\label{def:propensity}
\index{propensity score}
The propensity score is $e(x) = \Prob(A = 1 \mid X = x)$.
\end{definition}

The propensity score reduces a covariate vector of arbitrary dimension to a
scalar, and it does so without discarding any information relevant to
confounding. That claim is the content of the following result, which is due in this form
to \citet{marchetti2004balancing}.

\begin{theorem}[Balancing property]
\label{thm:balancing}
$X \indep A \mid e(X)$.
\end{theorem}

\begin{proof}
It suffices to show $\Prob\left( A = 1 \mid X, e(X) \right) = \Prob\left( A = 1 \mid e(X) \right)$.
The left side is $e(X)$ by Definition~\ref{def:propensity}, since $e(X)$ is a
function of $X$. For the right side, condition on $e(X)$ and apply the tower
property:
\begin{equation*}
  \Prob\left( A = 1 \mid e(X) \right)
    = \E\left[ \Prob\left( A = 1 \mid X \right) \,\middle|\, e(X) \right]
    = \E\left[ e(X) \mid e(X) \right]
    = e(X).
\end{equation*}
The two conditional probabilities agree almost surely, which is the assertion.
\end{proof}

\begin{corollary}[Ignorability given the score]
\label{cor:ps-ignorability}
Under Assumption~\ref{ass:ignorability},
$\left( Y(0), Y(1) \right) \indep A \mid e(X)$.
\end{corollary}

\begin{proof}
For $a \in \{0,1\}$ write
$\Prob\left( A = 1 \mid Y(a), e(X) \right) = \E\left[ \Prob\left( A = 1 \mid Y(a), X \right) \mid Y(a), e(X) \right]$.
Assumption~\ref{ass:ignorability} makes the inner probability equal to $e(X)$,
and the outer expectation of $e(X)$ given $e(X)$ is $e(X)$, which does not
depend on $Y(a)$.
\end{proof}

Corollary~\ref{cor:ps-ignorability} is what makes the score useful. Adjustment
for a scalar suffices where adjustment for a high dimensional covariate vector
was required, provided the score is known. It is not known, and the practical
difficulties in this chapter all descend from that.

\begin{remark}
\label{rem:estimated-score}
Estimating $e$ and then treating $\hat{e}$ as if it were known produces a
conservative standard error for the ATT under some conditions and an
anticonservative one under others. The efficient influence function
of Section~\ref{sec:dr-influence} accounts for the estimation of $e$ and should
be used in preference to a plug in variance.
\end{remark}

\section{Weighting estimators}
\label{sec:ps-weighting}

The inverse probability weighted estimator\index{inverse probability weighting}
of the ATE is
\begin{equation}
  \hat{\ATE}_{\mathrm{IPW}}
    = \frac{1}{n} \sum_{i=1}^{n}
      \left( \frac{A_i Y_i}{\hat{e}(X_i)}
           - \frac{(1 - A_i) Y_i}{1 - \hat{e}(X_i)} \right).
  \label{eq:ipw}
\end{equation}
Its unbiasedness under a correctly specified score follows from a one line
calculation:
\begin{equation}
  \E\left[ \frac{A Y}{e(X)} \right]
    = \E\left[ \frac{\E\left[ A \mid X \right] \E\left[ Y(1) \mid X \right]}{e(X)} \right]
    = \E\left[ \E\left[ Y(1) \mid X \right] \right]
    = \E\left[ Y(1) \right],
  \label{eq:ipw-unbiased}
\end{equation}
where the first equality uses Assumption~\ref{ass:ignorability} to factor the
conditional expectation and Assumption~\ref{ass:consistency} to
replace $Y$ by $Y(1)$ on the event $\{A = 1\}$.

The Hajek variant\index{Hajek estimator} normalises by the realised weights,
\begin{equation}
  \hat{\ATE}_{\mathrm{Hajek}}
    = \frac{\sum_i A_i Y_i / \hat{e}(X_i)}{\sum_i A_i / \hat{e}(X_i)}
    - \frac{\sum_i (1 - A_i) Y_i / (1 - \hat{e}(X_i))}{\sum_i (1 - A_i) / (1 - \hat{e}(X_i))},
  \label{eq:hajek}
\end{equation}
which is biased in finite samples and almost always has smaller variance. The
normalisation also makes the estimator invariant to a location shift in $Y$,
which~\eqref{eq:ipw} is not, and an estimator whose value depends on whether the
outcome is measured in earnings or in earnings minus a constant is not one to
report.

\section{Variance under weak overlap}
\label{sec:ps-variance}

Table~\ref{tab:overlap} showed 33 control units in the top propensity decile.
Their weights under~\eqref{eq:ipw} are of order $1/(1 - 0.96) = 25$, against a
mean control weight near 1.7. Table~\ref{tab:weighting-comparison} records what
this does to the estimates.

\begin{table}[t]
\centering
\caption[Weighting estimators on the retraining data]{Weighting estimators
applied to the retraining data. Effects are in units of annual earnings.
Standard errors are from 2000 bootstrap replications resampling units. The
trimmed rows discard units with $\hat{e}$ outside the stated interval and
therefore estimate a different, data defined estimand.}
\label{tab:weighting-comparison}
\begin{tabular}{lrrrr}
\toprule
Estimator & Estimate & Std. error & Max weight & Eff. $n$ \\
\midrule
Unadjusted difference        & 3184 & 141 & 1.0  & 8412 \\
IPW, Equation~\eqref{eq:ipw} & 1508 & 612 & 91.4 & 1147 \\
Hajek, Equation~\eqref{eq:hajek} & 1461 & 388 & 91.4 & 1147 \\
Hajek, trimmed to $[0.05, 0.95]$ & 1424 & 251 & 19.8 & 3908 \\
Hajek, trimmed to $[0.10, 0.90]$ & 1397 & 214 &  9.6 & 5216 \\
Overlap weights              & 1412 & 198 &  0.25 & 6841 \\
\bottomrule
\end{tabular}
\end{table}

Three readings of Table~\ref{tab:weighting-comparison} are worth separating. The
unadjusted contrast is more than twice the adjusted one, which is the expected
direction: units with better prior earnings enrolled more often. The untrimmed
IPW standard error is four times the trimmed one, and the effective sample size
falls to 14 percent of the data, which is the cost of the 33 units in the top
decile. And the point estimates across the adjusted rows agree to within 8
percent, which is reassuring about the adjustment and says nothing about the
unmeasured confounding of Figure~\ref{fig:confounding-dag}.

The overlap weights\index{overlap weights} of \citet{okafor2011overlap} in the
final row take the form
$w_i = 1 - \hat{e}(X_i)$ for treated units and $\hat{e}(X_i)$ for controls.
They target the effect in the subpopulation with the most equipoise, they are
bounded by construction, and they minimise the asymptotic variance among all
weighting schemes of the form $w(x) h(x)$. The price is that the estimand is
defined by the covariate distribution of the sample rather than chosen in
advance, which conflicts with Remark~\ref{rem:estimand-first} and should be
declared rather than discovered.

\section{Assessing balance}
\label{sec:ps-balance}

A fitted propensity model is checked by whether it balances covariates, not by
whether it predicts treatment well. A model with excellent discrimination
implies poor overlap and is bad news rather than good.

\begin{definition}[Standardised mean difference]
\label{def:smd}
\index{standardised mean difference}
For covariate $X_k$ with weighted arm means $\bar{x}_{k1}, \bar{x}_{k0}$ and
pooled unweighted variance $s_k^{2}$, the standardised mean difference is
$d_k = (\bar{x}_{k1} - \bar{x}_{k0}) / s_k$.
\end{definition}

Balance is conventionally declared at $|d_k| < 0.1$ for every covariate, a
threshold with no theoretical basis that has nevertheless proved a serviceable
default. What matters more than the threshold is that balance be assessed on
the same weights used for the estimate, including any trimming, and that
interactions and squared terms be checked rather than main effects alone. A
score that balances every main effect and no interaction is a score fitted to
the wrong model.
